\section{Local treatment geometry and the Gaussian transfer}

\subsection{Global tail control and local treatment geometry}

The recovery guarantee in \cref{obj:thm:adaptive-minimax} rests on global tail control and the equal-cell construction. To clarify the geometric scope of this guarantee, we compare the model with the local anti-concentration condition \leanref{sym:A3Dorn}{\ensuremath{\mathsf A_3^{\mathrm{Dorn}}}}. This condition relates a neighborhood's largest selected propensity to the covariate probability assigned to propensities of comparable size. Its formulation fixes both the family of laws and the pointwise propensity representatives.

\begin{definitionv}{def:dorn-a3}[Local anti-concentration condition $\mathsf A_3^{\mathrm{Dorn}}$]
The family-level condition
\(\mathsf A_3^{\mathrm{Dorn}}(\mathfrak P,(e_P)_{P\in\mathfrak P})\)
is defined for a family and selected propensity versions satisfying:
\begin{itemize}
\item \textbf{(Family.)} \(\mathfrak P\) is a nonempty family of probability laws with common covariate cube \(\mathcal S\).
\item \textbf{(Versions.)} For each \(P\in\mathfrak P\), \(e_P:\mathcal S\to[0,1]\) is a fixed, pointwise-defined measurable version of \(P(A=1\mid X=\cdot)\).
\item \textbf{(Overlap.)} For each \(P\in\mathfrak P\), \(e_P(X)\in(0,1)\) \(P\)-almost surely.
\end{itemize}
The condition holds when there exist constants \(\rho,\nu>0\) and \(h_0>0\), common to the family, such that, for every \(P\in\mathfrak P\), every \(x_0\in\mathcal S\), and every \(0<h\le h_0\),
\[
P_P\left(
e_P(X)\ge
\rho\sup_{\substack{x\in\mathcal S\\\|x-x_0\|_2\le h}}e_P(x)
\,\middle|\,
\|X-x_0\|_2\le h
\right)>\nu.
\]
Here \(P_P\) denotes probability under \(P\), and the supremum is the pointwise supremum of the selected version. Thus the condition is a property of the family together with its selected propensity versions; its truth can depend on their values on \(P_X\)-null sets. The covariate cube is \(\mathcal S=[-1,1]^d\) in Dorn's source and \(\mathcal S=\mathcal X\) after affine transport. For a single displayed law and propensity, \(\mathsf A_3^{\mathrm{Dorn}}(P)\) abbreviates the condition for the singleton family equipped with that selected version.
\end{definitionv}

The common constants in \cref{obj:def:dorn-a3} require this comparison to persist across laws, locations, and sufficiently small neighborhoods. The pointwise supremum also makes the selected representative part of the condition. By contrast, the global tail restriction in \cref{obj:ass:global-tail} controls the probability of small propensity values under the covariate distribution. A thin-strip construction makes the distinction concrete: it combines a radial propensity component with additional treatment probability concentrated in a narrowing strip.

\begin{definitionv}{def:thin-strip}[Thin-strip law $P_{\mathrm{strip}}$]
The observed-data law \(P_{\mathrm{strip}}\) is defined by the following construction, under these conditions:
\begin{itemize}
\item \textbf{(Dimension.)} \(d\ge2\).
\item \textbf{(Covariates.)} \(X\) is uniform on \(\mathcal X\).
\item \textbf{(Strip exponent.)} Choose \(a\) satisfying
\[
1<a<1+\frac{d}{\gamma-1}.
\]
\end{itemize}
Define the cube center, radial distance, and its continuous distribution function by
\[
x^\circ=(1/2,\ldots,1/2),
\qquad
R(x)=\|x-x^\circ\|_\infty,
\qquad
F_R(t)=P(R(X)\le t).
\]
Define the strip and selected propensity by
\[
S=\left\{x\in\mathcal X:
|x_2-x_2^\circ|\le |x_1-x_1^\circ|^a\right\},
\qquad
e_\star(x)=
\min\left\{1,F_R(R(x))^{1/(\gamma-1)}
+\frac14\mathbf 1_S(x)\right\}.
\]
Given \(X\), draw
\[
A\sim\operatorname{Bernoulli}(e_\star(X)).
\]
Set
\[
p_0=\min\left\{\frac14,\frac{L}{4B}\right\}.
\]
Draw \(Y(0)\) and \(Y(1)\) conditionally independently given \(X\), each with conditional distribution \(B\operatorname{Bernoulli}(p_0)\), and define
\[
Y=AY(1)+(1-A)Y(0).
\]
Then \(P_{\mathrm{strip}}\) is the law of \((X,A,Y)\).
\end{definitionv}

The thin-strip law \leanref{sym:Pstrip}{\ensuremath{P_{\mathrm{strip}}}} of \cref{obj:def:thin-strip} has constant potential-outcome means and bounded outcomes. It therefore isolates the treatment geometry while keeping the response smooth. The strip exponent makes the strip narrow relative to the surrounding neighborhood as its radius decreases. Under the displayed exponent range, \cref{obj:prop:strict-enlargement} establishes the coexistence of global tail control and a vanishing normalized treated-design moment.

\begin{propositionv}{prop:strict-enlargement}[Thin-strip enlargement and design degeneration]
Suppose the parameters satisfy:
\begin{itemize}
\item \textbf{(Dimension and smoothness.)} \(d\in\mathbb N\), \(d\ge2\), and \(\beta>1\).
\item \textbf{(Model constants.)} \(B>0\), \(L>0\), \(C\ge1\), and \(0<c_f\le1\).
\item \textbf{(Overlap and strip exponent.)} \(\gamma>1\) and
\[
1<a<1+\frac{d}{\gamma-1}.
\]
\end{itemize}
For every such \(a\), let \(P_{\mathrm{strip}}\) and its displayed selected propensity \(e_\star\) be the construction in \cref{obj:def:thin-strip}, and let \(\mathcal P_\gamma\) denote the model class with the fixed parameters \((d,\beta,B,L,C,c_f,\gamma)\). Then \(P_{\mathrm{strip}}\in\mathcal P_\gamma\), and
\[
P_{\mathrm{strip}}\{e_\star(X)\le t\}\le t^{\gamma-1},
\qquad t\in[0,1].
\]

On the unit cube \(\mathcal X=[0,1]^d\), the numerical anti-concentration clause fails: there are no \(\rho,\nu,h_0>0\) such that, for every \(x_0\in\mathcal X\) and every \(0<h\le h_0\),
\[
\frac{
P_{\mathrm{strip}}\!\left\{
e_\star(X)\ge
\rho\sup_{\substack{x\in\mathcal X\\ \|x-x_0\|_2\le h}}e_\star(x),
\ \|X-x_0\|_2\le h
\right\}
}{
P_{\mathrm{strip}}\{\|X-x_0\|_2\le h\}
}
>\nu.
\]
This failure applies to the numerical clause allowing propensity values equal to one. Consequently, the full singleton-family predicate
\(\mathsf A_{3}^{\mathrm{Dorn}}(P_{\mathrm{strip}})\), evaluated using \(e_\star\), also fails.

Moreover, writing \(x^\circ=(1/2,\ldots,1/2)\), the normalized treated-design second moment in the second coordinate satisfies
\[
\lim_{h\downarrow0}
\frac{
\displaystyle
\int_{\{A=1,\ \|X-x^\circ\|_\infty\le h\}}
\left(\frac{X_2-1/2}{h}\right)^2
\,dP_{\mathrm{strip}}
}{
P_{\mathrm{strip}}\{A=1,\ \|X-x^\circ\|_\infty\le h\}
}
=0.
\]
In particular, \(\mathcal P_\gamma\) contains a law whose full singleton-family Dorn predicate fails for the displayed selected propensity \(e_\star\).
\end{propositionv}

The global tail inequality controls the prevalence of small propensities, while the displayed moment records the shape of the treated design near the cube center. Its limit shows that, after normalization by neighborhood width, treated observations concentrate increasingly close to the strip's central hyperplane. The local second-coordinate variation consequently degenerates even though the law satisfies the global model restrictions. For dimensions at least two, \cref{obj:prop:strict-enlargement} thus establishes a strict enlargement relative to imposing the selected-version predicate within the causal class. The equal-cell estimator's guarantee in \cref{obj:thm:adaptive-minimax} covers this geometry through its template and count-feasibility conditions.

\subsection{Transfer to fixed-constant Gaussian source envelopes}

The geometric distinction also has a source-model consequence. Under the retained restrictions in \citet[Section 2, Assumptions 1--2, pp.~3--7]{Dorn2026GlobalOverlap}, affine transport maps each fixed-constant Gaussian source envelope to the causal class used by \cref{obj:thm:adaptive-minimax}. A common response bound controls the transported mean and noise scale, while fixed-cube smoothness completion supplies a common full Hölder radius. Consequently, one transformed equal-cell estimator attains the expected spatial-supremum rate uniformly over the source envelope and the prescribed compact exponent range. The exact envelope, transport constants, exponent-specific Gaussian witness, and family-richness qualifications are stated in \cref{obj:prop:dorn-a3-free-corollary} and proved in the source-comparison appendix.



The transfer aligns the source restrictions with the components needed for response recovery. Uniform covariates supply the density lower bound after transport, the propensity tail inequality supplies global overlap control, and the Gaussian and treated-moment restrictions supply common response bounds. Fixed-cube smoothness completion converts the source's top-derivative seminorm restriction, together with a function bound, into the full Hölder radius used by the estimator. The precise completion argument is given in \cref{obj:lem:fixed-cube-holder-completion}.

The upper guarantee controls the expected spatial supremum over the entire unit cube, uniformly over the fixed-constant envelopes and the compact exponent range. This loss criterion is stronger than the displayed essential-supremum probability criterion: it controls the magnitude of the largest spatial error in expectation and yields the stated probability guarantee at every positive tolerance. The same constants apply to each source subfamily satisfying the retained restrictions because each such family lies within the corresponding envelope.

The Gaussian existence assertion has a separate quantifier structure. In dimensions at least two, it supplies, for each overlap exponent, a source tuple satisfying the retained restrictions with positive finite constants chosen for that exponent and with local anti-concentration failure for its selected representative. These witnesses establish the geometric scope of the retained Gaussian restrictions; the envelope-wide upper guarantee applies with the fixed constants specified at the start of \cref{obj:prop:dorn-a3-free-corollary}.

The converse statements likewise identify their comparison classes explicitly. A source family with a common treated curve has zero minimax pointwise miss probability, whereas the full causal class in \cref{obj:def:risks} admits the displayed positive expected absolute-error lower bound at the oracle scale, supported by its two-valued potential-outcome subclass. Family richness therefore determines the scope of a minimax converse, while the transferred upper guarantee remains valid for every qualifying source subfamily. The technical source-family comparison is developed in \cref{obj:lem:arbitrary-family-lower-obstruction,obj:lem:dorn-rate-scope}.
